\documentclass[11pt]{article}
\usepackage[margin=1in]{geometry}
\usepackage{amsmath,amssymb,booktabs,graphicx}
\usepackage[colorlinks=true,linkcolor=teal]{hyperref}
\title{A Small Guide to Better Experiments}
\author{Research notes · LaTeX Coder}
\date{September 2026}
\begin{document}
\maketitle

\begin{abstract}
Good experiments turn a question into evidence. This working paper combines mathematical notation, structured observations, and reproducible results in one editable document.
\end{abstract}

\section{Start with a clear question}
We study how \textbf{sample size} affects the stability of an estimate. The goal is simple: make each decision \emph{explicit}, keep the analysis reproducible, and communicate uncertainty alongside the result.

For independent observations $x_1,\ldots,x_n$, the sample mean is $\bar{x}=\frac{1}{n}\sum_{i=1}^{n}x_i$. The uncertainty shrinks as more observations become available.

\subsection{The estimation model}
\begin{equation}
\mathcal{L}(\theta)=\frac{1}{n}\sum_{i=1}^{n}\left(y_i-f_\theta(x_i)\right)^2+\lambda\lVert\theta\rVert_2^2
\end{equation}

The regularization parameter $\lambda\geq 0$ balances goodness of fit and model complexity. We update the parameters using the gradient:

\[
\theta_{t+1}=\theta_t-\eta\nabla_\theta\mathcal{L}(\theta_t),\qquad \eta=0.01.
\]

\section{Record the evidence}
The pilot compares three sample sizes under identical conditions. The larger run reduces error while keeping the confidence interval easy to interpret.

\begin{table}[ht]
\centering
\begin{tabular}{lrrr}
\toprule
\textbf{Experiment} & \textbf{Samples} & \textbf{Mean error} & \textbf{Time (s)} \\
\midrule
Pilot & 32 & 0.184 & 1.2 \\
Standard & 128 & 0.092 & 3.8 \\
Extended & 512 & 0.046 & 12.4 \\
\bottomrule
\end{tabular}
\caption{A controlled comparison of accuracy and computation time.}
\end{table}

\subsection{What changed?}
\begin{itemize}
\item \textbf{Accuracy improved:} the extended run cuts mean error by 75\%.
\item \textbf{Cost increased:} more observations require additional computation.
\item \textbf{Uncertainty matters:} report the interval, not just the point estimate.
\end{itemize}

\section{Make the next run reproducible}
\begin{enumerate}
\item Fix the random seed and record the data version.
\item Repeat each experiment with five independent trials.
\item Review outliers before aggregating the results.
\end{enumerate}

\begin{align}
\mathbb{E}[\bar{x}] &= \mu, \\
\operatorname{Var}(\bar{x}) &= \frac{\sigma^2}{n}.
\end{align}

\section{A practical conclusion}
Write the explanation next to the evidence. A useful paper lets readers follow the argument from the \emph{question} to the \textbf{decision}, with every assumption visible.

\begin{quote}
Clarity is a property of the whole document: prose, equations, tables, and the source that connects them.
\end{quote}
\end{document}
